\documentclass{plantillaPPt}
\titulo{14}{Series de Taylor y Fin de Cálculo II}
\begin{document}
\primeraDiapo

\begin{frame}{Problema 1}
    \framesubtitle{Práctica 14}
    1. Verifique que la serie de potencias de $f(x)=(4+x^2)^{-1}$ es
    \begin{equation*}
        \sum_{n=0}^\infty (-1)^n\frac{x^{2n}}{4^{(n+1)}}
    \end{equation*}
    indicando su radio de convergencia.
\end{frame}

\begin{frame}{Polinomio de Taylor}
\framesubtitle{Definicion}

\begin{definicion}
    Sea \(f\) una función con derivadas de orden \(n\) en un punto \(a\). El \textbf{Polinomio de Taylor de orden \(n\)} generado por \(f\) en \(a\) es:
    \[
    P_n(x) = f(a) + f'(a)(x-a) + \frac{f''(a)}{2!}(x-a)^2 + \cdots + \frac{f^{(n)}(a)}{n!}(x-a)^n
    \]
    En notación de sumatoria:
    \[ P_n(x) = \sum_{k=0}^n \frac{f^{(k)}(a)}{k!} (x-a)^k \]
\end{definicion}

\end{frame}

\begin{frame}{Teorema de Taylor (Resto de Lagrange)}
\framesubtitle{Estimación de errores}

\begin{teorema}
    Si \(f\) es derivable \(n+1\) veces en un intervalo \(I\) que contiene a \(a\), entonces para cada \(x \in I\), existe un número \(c\) entre \(a\) y \(x\) tal que:
    \[ f(x) = P_n(x) + R_n(x) \]
    Donde \(R_n(x)\) es el \textbf{Resto de Lagrange}:
    \[ R_n(x) = \frac{f^{(n+1)}(c)}{(n+1)!}(x-a)^{n+1} \]
\end{teorema}
\vspace{-10pt}
Esto nos permite medir el error de la aproximación: \(|f(x)-P_n(x)| = |R_n(x)|\).

\end{frame}

\begin{frame}{Problema 2}
    \framesubtitle{Práctica 14}
    2. Sea $f$ la función definida por $f(x) = \cos(3x+\frac{\pi}{6})$ y $P(x)$ el polinomio de Taylor de grado 4 de $f$ con centro $x=0$.\\[10pt]
    $(a)$ Determianr $P(x)$.\\
    $(b)$ Utilizar el resto de Lagrange para mostrar que 
    \begin{equation*}
        \left|f\left(\frac{1}{6}\right)-P\left(\frac{1}{6}\right)\right|<\frac{1}{3000}.
    \end{equation*}
\end{frame}
\begin{frame}{Problema 3}
    \framesubtitle{Práctica 14}
    3. La serie de Taylor de $f$ centrada en $x=5$ converge a $f(x)$ para todo $x$ en su intervalo de convergencia. La $n-$ésima derivada de $f$ en $x=5$ está dada por
    \begin{equation*}
        f^{(n)}(5) = \frac{(-1)^nn!}{2^n(n+2)} \text{  con  }f(5)=\frac{1}{2}.
    \end{equation*}
    $(a)$ Determinar el polinomio de Taylor de grado 3 de $f$ centrado en $x=5$.\\
    $(b)$ Mostrar que el polinomio de Taylor del item $(a)$ aproxima $f(6)$ con un error menor a $0.02$.
\end{frame}
\begin{frame}{Problema 4}
    \framesubtitle{Práctica 14}
    4. Sabiendo que $\displaystyle \forall x\in\mathbb{R}, e^x=\sum_{n=0}^\infty \frac{x^n}{n!}$, expresar $f(x) = e^{-x^2}$ como una serie y determinar el valor de la integral
    \begin{equation*}
        \int_0^1 e^{-x^2}\,dx
    \end{equation*}
    con un error menor a $\frac{1}{200}$.
\end{frame}
\end{document}